{\tiny
\setlength{\tabcolsep}{3pt}
\begin{longtable}{>{\raggedright\arraybackslash}p{2.55cm}>{\raggedright\arraybackslash}p{2.55cm}>{\raggedright\arraybackslash}p{2.55cm}>{\raggedright\arraybackslash}p{2.55cm}>{\raggedright\arraybackslash}p{2.55cm}>{\raggedright\arraybackslash}p{2.55cm}}
\caption{Registre des paramètres}\label{tab:registre_parametres}\\
\toprule
\textbf{Paramètre} & \textbf{Valeur} & \textbf{Statut} & \textbf{Justification} & \textbf{Références} & \textbf{Calibrage} \\
\midrule
\endfirsthead
\toprule
\textbf{Paramètre} & \textbf{Valeur} & \textbf{Statut} & \textbf{Justification} & \textbf{Références} & \textbf{Calibrage} \\
\midrule
\endhead
\bottomrule
\endfoot
tol\_rel & 0,000001 & benchmark & Plus grande tolérance relative pour laquelle les résultats (critère final, $\hat K$, ARI) ne diffèrent pas de ceux obtenus avec une tolérance 100 fois plus stricte, au-delà de l'écart observé entre deux relances convergeant vers le même optimum. Grille 1e-3 à 1e-10 sur les jeux du protocole. &  & tolerance \\
n\_iter\_max & 500 & benchmark & Plafond d'étapes M tel que, avec tol\_rel calibrée, moins de 1 \% des ajustements du pilote l'atteignent ; les ajustements plafonnés sont marqués non convergés. &  & tolerance \\
reg\_covariance & 0,000001 & recherche & Régularisation + reg I des covariances (EM, MAP-EM). Choisie sur une grille logarithmique 0 et 1e-10 à 1e-1 en maximisant la log-vraisemblance moyenne hors échantillon (validation croisée sur les jeux d'apprentissage). &  & regularisation \\
kmeans\_n\_iter\_max & 100 & benchmark & Plafond des itérations de Lloyd pour l'initialisation : la part des initialisations qui l'atteignent est mesurée sur les jeux du protocole ; le plafond est relevé (par pas de 100) tant qu'elle dépasse 1 \%. &  & kmeans \\
n\_init & {'synthetique': 10, 'iris': 10, 'sp500': 10, 'climat': 10, 'texte': 10} & benchmark & Par famille de données : nombre de relances tel que la probabilité qu'au moins une relance atteigne le meilleur optimum (à 10 fois la tolérance d'arrêt près) soit au moins 0,99, soit ceil(ln 0,01 / ln(1 - p)) où p est estimé par 200 relances pilotes, pour le plus difficile des algorithmes comparés ; plafonné à la taille du pilote (200), la probabilité résiduelle étant alors rapportée. & biernacki2003 & n\_init \\
seuil\_actif & 0,010000 & plan & Composante active si N\_k / N >= 1 \% (plan de recherche) ; sensibilité aux seuils 0,5 \%, 5 \% et N\_k >= 1 rapportée (scénario S10). &  &  \\
k\_max & {'synthetique': 6, 'climat': 10, 'sp500': 6} & plan & Nombre maximal de composantes du plan (au moins deux fois le K de référence) ; la dépendance à Kmax est un objet d'étude (Iris, sections 2.3.11), non un réglage. &  &  \\
alpha0\_map & 1,000000 & mathematique & MAP-EM : plus petite concentration de Dirichlet pour laquelle le mode existe à l'intérieur du simplexe (alpha0 >= 1) ; alpha0 = 1 n'ajoute aucun pseudo-compte, le mode en pi coïncide avec l'estimateur du maximum de vraisemblance. &  &  \\
priors\_reference & emp,diff & source & Deux jeux de référence comparés tels quels : « empiriques » (beta0 = 1, nu0 = D, W0\^{}-1 = covariance empirique, m0 = moyenne, alpha0 = 1/Kmax : valeurs par défaut documentées de BayesianGaussianMixture) et « diffus » (beta0 = 1e-3, W0 = I, nu0 = D). Ils ne sont pas présentés comme optimaux : le bras calibré est VB-cal. & pedregosa2011 &  \\
espace\_priors & {'alpha0': (0.001, 10.0), 'beta0': (0.0001, 100.0), 'nu0\_plus': (0.0, 50.0), 'echelle\_w0': (0.01, 100.0)} & plan & Espace de recherche des hyperparamètres Normal-Wishart et Dirichlet (bornes, échelle logarithmique sauf nu0 - D) : il contient les deux jeux de référence (alpha0 = 1/Kmax, beta0 = 1e-3 ou 1, nu0 = D, W0 = I ou covariance empirique) et s'étend de deux ordres de grandeur au moins de part et d'autre ; nu0 > D - 1 garantit une loi de Wishart propre (Bishop 2006, section 2.3.6). Les illustrations de sensibilité parcourent les mêmes bornes. & bishop2006 &  \\
priors\_vb\_calibres & {'synthetique': {'alpha0': 0.1, 'beta0': 1.0, 'nu0\_plus': 0.0, 'echelle\_w0': 1.0}, 'climat': {'alpha0': 0.1, 'beta0': 1.0, 'nu0\_plus': 0.0, 'echelle\_w0': 1.0}, 'sp500': {'alpha0': 0.1, 'beta0': 1.0, 'nu0\_plus': 0.0, 'echelle\_w0': 1.0}} & recherche & Bras VB-cal et MAP-cal : hyperparamètres choisis par recherche aléatoire puis divisions successives (successive halving) sur l'espace « espace\_priors », en maximisant la log-vraisemblance prédictive moyenne en validation croisée sur les données d'apprentissage, séparément pour chaque famille de données (plis spatiaux par blocs pour le climat). Toutes les évaluations sont conservées pour l'analyse de sensibilité (configurations plausibles). & bergstra2012, jamieson2016 & priors \\
alpha\_test & 0,050000 & plan & Niveau des tests, déclaré dans le plan ; contrôle du risque familial par la procédure de Holm. & holm1979 &  \\
marges\_equivalence & {'ARI': 0.03, 'LL': 0.02} & donnees & Plus petite différence d'intérêt (Lakens 2017) : quantile à 95 \% de l'écart absolu entre deux ajustements de la MÊME méthode avec des graines différentes sur les mêmes données (plancher de bruit d'optimisation), avec un plancher de 0,001 ; un écart inférieur n'est pas attribuable à la méthode. & lakens2017, schuirmann1987 & marges \\
marge\_k & 0,500000 & mathematique & Marge d'équivalence sur l'erreur moyenne de K : K est entier, un écart moyen inférieur à 0,5 signifie que la majorité des réplications ne diffèrent pas d'une composante entière. &  &  \\
n\_bootstrap\_ic & 2000 & benchmark & Nombre de rééchantillonnages des intervalles bootstrap : erreur Monte Carlo des bornes inférieure à la moitié de la dernière décimale affichée (3 décimales), estimée par répétition du bootstrap. &  & monte\_carlo \\
n\_bootstrap\_modele & 1000 & benchmark & Nombre de réplicats bootstrap des intervalles de paramètres et de VaR : même règle de précision Monte Carlo. &  & monte\_carlo \\
n\_tirages\_posterieur & 4000 & benchmark & Tirages a posteriori pour les intervalles de crédibilité : même règle de précision Monte Carlo. &  & monte\_carlo \\
n\_simulation\_var & 200000 & benchmark & Tirages Monte Carlo de la VaR et de l'ES prédictives : même règle de précision Monte Carlo. &  & monte\_carlo \\
effet\_standardise & 0,500000 & source & Hypothèses de différence sans marge naturelle (taux d'ajustements dégénérés, écart relatif de perplexité) : la puissance est calculée pour un écart vrai égal à 0,5 écart-type des différences appariées du pilote (effet « moyen » de Cohen) ; pour les métriques munies d'une marge (ARI, log-vraisemblance, K), l'écart vrai est la marge. & cohen1988 &  \\
taille\_pilote & 10 & plan & Réplications pilotes servant à estimer la variance des différences appariées avant le calcul de puissance ; les données du pilote ne sont pas réutilisées dans les tests confirmatoires (graines disjointes). & lakens2017 &  \\
replications & {'scenarios': 30, 'climat\_sous': 40, 'climat\_semi': 60, 'climat\_d20': 40, 'climat\_sensibilite': 20, 'texte': 10} & benchmark & Nombre de réplications : puissance d'au moins 0,8 du test déclaré (TOST à la marge calibrée, ou Wilcoxon pour un écart égal à la marge) au niveau corrigé, d'après la variance des différences observée dans le pilote (Lakens 2017). & lakens2017 & puissance \\
repetitions\_chronometrage & 20 & plan & Répétitions de chaque mesure de temps (un seul processus, un seul fil) : avec 20 mesures, l'intervalle entre les 6e et 15e statistiques d'ordre contient la médiane avec une probabilité 1 - 2 P(Bin(20 ; 0,5) <= 5) = 0,959, sans hypothèse de loi ; l'intervalle bootstrap de la médiane du rapport des temps est rapporté. &  &  \\
taille\_bloc\_km & 100,000000 & donnees & Côté des blocs de la validation croisée spatiale : au moins la portée de l'autocorrélation spatiale (Roberts et al. 2017). Portée pratique = plus petite distance où le variogramme empirique exhaustif des variables standardisées atteint 95 \% de son palier ; pour des variables réduites, E||x\_i - x\_j||² / (2D) = 1 pour deux points indépendants, le palier vaut donc 1. & roberts2017 & spatial \\
n\_plis & 5 & plan & Nombre de plis de la validation croisée par blocs (plan) ; l'affectation des blocs aux plis est tirée au hasard. &  &  \\
espace\_lda & {'alpha': (0.001, 10.0), 'lissage': (0.0001, 1.0)} & plan & Espace de recherche de la LDA (échelle logarithmique) : concentration de Dirichlet sur theta\_d sur les mêmes bornes que alpha0 (« espace\_priors ») ; lissage des distributions de mots d'un pseudo-compte négligeable (1e-4) à un pseudo-compte par mot et par thème (1). & blei2003 &  \\
lda\_hyper & {'alpha': 0.1, 'lissage': 0.001} & recherche & Hyperparamètres de la LDA : recherche aléatoire puis divisions successives sur « espace\_lda », en minimisant la perplexité par complétion de documents de validation tirés du réservoir d'apprentissage (jamais des documents de test). & bergstra2012, jamieson2016 & texte \\
vocabulaire & {'taille': 1500, 'df\_min': 5, 'df\_max': 0.5, 'longueur\_min\_mot': 3, 'longueur\_min\_doc': 20, 'sensibilite': (500, 3000)} & plan & Prétraitement fixé avant l'analyse et identique pour pLSA et LDA : mots présents dans au moins 5 documents et au plus la moitié d'entre eux (les mots vides sont écartés par cette fréquence documentaire, sans liste externe), mots d'au moins 3 lettres après racinisation de Porter (1980), documents d'au moins 20 mots retenus. Ce n'est pas un réglage optimisé : la sensibilité des conclusions à la taille du vocabulaire (500 et 3000 mots) est mesurée. & porter1980 &  \\
espace\_variantes & {'kappa': (0.5, 1.0), 't0': (1.0, 1000.0), 'lot': (10, 150), 'epoques': (1, 50), 'sem\_n\_iter': (50, 1000), 'sem\_part\_burn': (0.0, 0.9), 'tau': (1.0, 1000.0)} & plan & Espace de recherche des variantes stochastiques : kappa dans ]0,5 ; 1] (conditions de Robbins et Monro sur les pas) ; décalages t0 et tau, tailles de lot (de 10 à la moitié du jeu de 300 points) et nombres d'époques en échelle logarithmique ; SEM : nombre d'itérations et part de chauffe. & robbins1951, cappe2009, hoffman2013 &  \\
variantes\_stochastiques & {'sem\_burn': 50, 'sem\_n\_iter': 150, 'online\_kappa': 0.6, 'online\_t0': 10.0, 'online\_lot': 30, 'svi\_kappa': 0.7, 'svi\_tau': 1.0, 'svi\_lot': 30, 'epoques': 15} & recherche & Pas et tailles de lot de la SVI et de l'EM online, burn-in du SEM : recherche sur la log-vraisemblance de validation, kappa contraint à ]0,5 ; 1] (conditions de Robbins et Monro). & robbins1951, cappe2009, hoffman2013 & variantes \\
hmm\_etats & 2 & benchmark & Nombre d'états du HMM choisi par BIC sur la période d'apprentissage (S de 1 à 5). & schwarz1978 & hmm \\
niveaux\_var & (0.05, 0.01) & plan & Niveaux de VaR du plan (1 \% et 5 \%). &  &  \\
elements\_lot\_max & 30000000 & performance & Taille maximale d'un tenseur (b, K, N) traité d'un coup : borne mémoire, sans effet sur les résultats (tests d'équivalence). &  &  \\
\end{longtable}
}
